\subsection{A closed theorem for the Case~I host family}
\label{sec:host-case-one}

\paragraph{K17-HOST-CASE-I: a theorem, beyond the field definitions.}
There is a closed source proof of the host-family assertion
\texttt{CountI}.  In full namespace notation its proof term is the
composition
\[
 \texttt{UpperTheory.V.U3.countI}\;
 \texttt{UpperTheory.V.SmallClassNumbers.corGK}.
\]
The first theorem has type
\(\texttt{U3.CorGK}\longrightarrow\texttt{HV.CountI}\);
the second is the closed theorem already used for the rational-field
specialization.  This composition uses the established, field-dependent
polynomial theorem \texttt{CorGK}.  It does not require the stronger
uniform discriminant and level exponents appearing in the separately
defined proposition \texttt{HV.CorG}.

Here is the precise mathematical reading.  Fix a character \(\chi\) in
the source's \texttt{PChar} domain, of modulus \(m\), and put
\(o=\operatorname{ord}(\chi)\).
Let \(p_K\) be the least prime at least seven congruent to three modulo
four and coprime to \(m\); set
\[
 K_\chi=\mathbb Q(\sqrt{-p_K}),\qquad
 h_\chi=|\operatorname{Cl}(K_\chi)|.
\]
Let \(p_1\) be the least prime congruent to one modulo the odd part of
\(o\), and coprime to \(2mp_K\).  The finite set used by the source is
\[
 \mathcal B_\chi=\{3\}\ \cup\
 \{\ell:\ell\text{ is an odd prime dividing }omp_Kh_\chi p_1\}.
\]
There exist natural numbers \(C\) and distinct primes \(r_1,r_2\),
depending only on \(\chi\), with the following property.

Choose a host \(h=(\chi,a,q,d)\), where \(a\ge2\), \(q,d\) are primes,
\(q\equiv3\pmod8\), \(d\equiv7\pmod{16}\), and
\(\gcd(qd,m)=1\).  These are exactly the source's
\texttt{PrimePair} conditions.  Its level and associated field are
\[
 N=\operatorname{lcm}(2^{a+2},m)qd,
 \qquad
 E_h=F_\chi\,\mathbb Q(\mu_{2^{a+2}})^+
          \mathbb Q(\sqrt{-q})\,\mathbb Q(\sqrt{-d}).
\]
Require
\[
 r_i<q,\quad r_i<d\quad(i=1,2),
\]
and, for \(t\in\mathbb N\), the explicit host-size conditions
\[
 2t+1<q,d,
 \qquad y<q,d\quad
 (y\in\{m,o,p_K,h_\chi,p_1\}).
\]
These are precisely \(h.\texttt{Big}(t)\).  Let \(e\in\mathbb N\) with \(0<e<t\), and let
\(L\) be any finite set of odd primes such that
\[
 \ell\notin\mathcal B_\chi,
 \qquad (\ell-1)\nmid 2e
 \qquad(\ell\in L).
\]
This pair of conditions is exactly \(h.\texttt{CaseI}(\ell,e)\).

For each \(\ell\), take any coefficient ring \(\mathcal O_\ell\) of the
source, and let \(X_h(\mathcal O_\ell,e)\) be the characters modulo
\(N\) that factor through \(\operatorname{Gal}(E_h/\mathbb Q)\) and
satisfy \(\psi(-1)=(-1)^e\).  Use \(V_{\ell,N}(e,\psi)\) for the
cohomology group with the everywhere-inertia and selected decomposition
conditions defined in the source reconstruction above.  Then all these
groups, for \(\ell\in L\) and \(\psi\in X_h(\mathcal O_\ell,e)\),
have finite length, and
\begin{equation}\label{eq:host-case-one-bound}
 \sum_{\ell\in L}\log\ell\!
 \sum_{\psi\in X_h(\mathcal O_\ell,e)}
 \frac{\operatorname{length}_{\mathcal O_\ell}
                   V_{\ell,N}(e,\psi)}
      {\operatorname{length}_{\mathcal O_\ell}
                   (\mathcal O_\ell/\ell\mathcal O_\ell)}
 \le (Nt)^C.
\end{equation}
The constants precede the host, the twist, the coefficient choices and
the finite prime set.  Both even and odd \(e\) are permitted here;
the character parity changes with \(e\).  No assertion about the whole
ambient \(H^1\) is substituted for these local kernels.

\paragraph{Proof and source check.}
Apply \texttt{CorGK} to the fixed imaginary quadratic field
\(K_\chi\), obtaining constants \(c_5,A\).  Choose
\(r_1\ge30mp_K+1\) and \(r_2\ge r_1+1\), both prime.
The prime divisors of the host level are among those dividing \(2mqd\).
The displayed inequalities therefore ensure
\(r_i\nmid30N|\operatorname{disc}(K_\chi)|\).
The source theorem \texttt{gdOfHost}, with its complete proof checked,
shows that every character in \(X_h(\mathcal O_\ell,e)\) satisfies the
\texttt{GD} hypotheses for this fixed \(K_\chi\).  Its field exclusion
is proved by ramification at \(p_K\): this prime ramifies in
\(K_\chi\), while the cyclotomic and character fields occurring in the
opposing compositum are unramified there.  The class-number and character
order exclusions are exactly supplied by \(\ell\notin\mathcal B_\chi\).

Group the characters by the kernel \(H\) of their character on
\((\mathbb Z/N\mathbb Z)^\times\).  Each quotient by such a kernel is
cyclic.  Applying \texttt{CorGK} to each resulting character orbit gives
finiteness and a bound with orbit factor \(\varphi([G:H])\), where
\(G=(\mathbb Z/N\mathbb Z)^\times\).  The source's finite character
count proves
\[
 \sum_H\varphi([G:H])\le |G|=\varphi(N)\le N.
\]
Indeed the characters with a fixed kernel are the faithful characters
of that cyclic quotient, numbering \(\varphi([G:H])\), and different
kernels give disjoint sets of characters.  Over a characteristic-zero
algebraically closed coefficient field there are exactly \(|G|\)
characters of the finite abelian group \(G\).  This also proves the
inequality without assuming the individual rings \(\mathcal O_\ell\)
already contain every root of unity.

Consequently the left side of \eqref{eq:host-case-one-bound} is at most
\[
 2c_5(\log r_1+\log r_2)\,N\,[N(e+1)]^A.
\]
Choose \(k\in\mathbb N\) with
\(2c_5(\log r_1+\log r_2)\le16^k\).
Since \(N\ge16\), \(e+1\le t\), and \(t\ge1\), this is at most
\((Nt)^{k+1+A}\).  Set \(C=k+1+A\).
These are the choices and inequalities in the checked source proof.

The exact implementation locators, relative to
\path{src/Development/AutomorphicArithmetic/AutomorphicVarieties/}, are
\path{UnitaryThreeGroups/CountI.lean}, theorem
\texttt{UpperTheory.V.U3.countI};
\path{SmallClassNumbers/CorGKOf.lean}, theorem
\texttt{UpperTheory.V.SmallClassNumbers.corGK};
\path{HeightAndVanishing/GdOfHost.lean}, theorem \texttt{gdOfHost};
\path{HeightAndVanishing/Counts.lean}, definition \texttt{CountI};
and \path{HeightAndVanishing/DirichletPrimeToLPartMulLPowerPart.lean},
definitions \texttt{Host}, \texttt{Big}, \texttt{X}, and \texttt{CaseI}.
The auxiliary primes and bad set are defined in \path{StmtFields.lean}.
The complete \texttt{U3.countI} proof, including its quantifier order,
auxiliary-prime exclusions, finite orbit partition and final absorption
inequality, was independently read.  No fresh Lean compilation is
claimed.

\paragraph{K17-HOST-POINT: the full point assertion also closes.}
Further retrieval supplies the remaining inputs to
\texttt{pointHV\_of\_counts}.  Consequently \texttt{PointHV} has a
closed proof by composition of the recovered source theorems.  Its
arithmetic layer-bound premise remains part of the theorem; closing
the source proof does not set that bound to zero.

For precision, fix \(\chi\) and any function
\(j:\mathbb N\longrightarrow\mathbb N\).  There are
\(C\in\mathbb N\) and primes \(r_1,r_2\), depending only on
\(\chi,j\), such that the following holds.  Take any host \(h\) as
above, with \(r_i<q,d\), any \(t\ge2\) satisfying
\(h.\texttt{Big}(t)\), and any real number \(B\ge0\).
For each \(\ell\in\mathcal B_\chi\), assume the following layer
condition: for every cyclotomic \(\ell\)-power extension
\[
 F/\mathbb Q(\mu_{\ell m}),\qquad
 [F:\mathbb Q(\mu_{\ell m})]=\ell^{j(\ell)},
\]
all the groups \(\mathcal S_\ell(F,e)\), \(1\le e<t\), have finite
length, and
\begin{equation}\label{eq:host-layer-premise}
 \sum_{e=1}^{t-1}\operatorname{length}_{\mathbb Z_\ell}
                 \mathcal S_\ell(F,e)\le B.
\end{equation}
Here \emph{cyclotomic \(\ell\)-power extension} means
\(F=\mathbb Q(\mu_{\ell m})(\mu_{\ell^b})\) for some \(b\), as
encoded by \texttt{IsLayer}; the statement is invariant under the
number-field presentations used in the source.

Write
\[
 \mathcal S_\ell(E_h,e)^{+_e}
 =\{x\in\mathcal S_\ell(E_h,e):c_{E_h}x=(-1)^e x\},
\]
where the complex-conjugation action includes the Tate coefficients.
Then these eigenspaces have finite length for every odd prime \(\ell\)
and every \(1\le e<t\).  For every finite set \(L\) of odd primes,
\begin{equation}\label{eq:host-point-bound}
 \sum_{\ell\in L}\log\ell\sum_{e=1}^{t-1}
 \operatorname{length}_{\mathbb Z_\ell}
       \mathcal S_\ell(E_h,e)^{+_e}
 \le (Nt)^C+CB.
\end{equation}
There is no Case~I restriction on \(L\) in this assertion.  This is
the precise field-level, finite-parameter statement defined by
\texttt{PointHV}, now with its source inputs supplied.

\paragraph{Closed source inputs and proof.}
Set \(h_N=\texttt{SmallClassNumbers.litN}\), and obtain
\(h_H\) by applying
\path{published_hypotheses_from_patching_alt} to \(h_N\) and the
closed source theorem \texttt{descWeak}.  The following checked
assembly supplies every argument of \texttt{pointHV\_of\_counts}.
The directory prefix in the next list is
\path{ModularityArgument/ThresholdArguments/DenseThresholdCases/}.
\begin{enumerate}
 \item \path{HVFinalU.lean}, \texttt{pt\_s}, gives Poitou--Tate
 duality from \(h_N\); the local compact finiteness and duality
 arguments are the fields \texttt{hN.coh.lfK} and
 \texttt{hN.coh.ltK}.
 \path{KFinal.lean}, \texttt{lemmaK\_litN\_alt}, gives
 \(\exists C_K,\ \texttt{LemmaK}(C_K)\).
 \item \texttt{CountI} is supplied by the closed composition proved
 above.  \texttt{CountII} and \texttt{CountIII} are supplied by
 \path{AdjointLiftingOfCohomology.lean},
 \path{countII_and_countIII_of_printsH_of_sp4In}, using
 \path{Sp4In0.lean}, \texttt{sp4In0\_alt}, at \(h_H\).
 \item \path{DnEndsTop.lean}, \texttt{ends0\_alt}, supplies the
 reduced unitary lifting statement at \(h_H\).  Applying
 \path{LinkedUnlE.lean},
 \path{count_v_of_printed_bound_k_reduced_unitary_lift}, with
 \(h_H\) and the preceding \texttt{LemmaK} witness gives
 \texttt{CountV}.
 \item Outside that directory, \path{LevelCohomology/PropZ.lean}
 proves \texttt{X.propZ} without additional arguments.
 \texttt{countZ X.propZ} gives \texttt{CountZ}.
\end{enumerate}
The construction of \(h_H\) is in
\path{PublishedHypothesesFromPatching.lean}; the exact same closed
\(h_N,h_H\) are used by \texttt{SmallClassNumbers.corGK}.
The theorem \texttt{pointHV\_of\_counts} is in
\path{HeightAndVanishing/JointWinOfCounts.lean}.  Applying it to the
inputs above proves \texttt{PointHV}.  This is a forward proof
composition; no converse implication from the later \texttt{VWin}
theorem is being used.

For clarity, the final arithmetic combination can also be checked
directly.  The source's \texttt{sumOut} bounds the contribution from
odd primes outside \(\mathcal B_\chi\) by \((Nt)^{C_0}\).
For each \(\ell\in\mathcal B_\chi\), its \texttt{sumIn}, using
\eqref{eq:host-layer-premise}, bounds that prime's contribution by
\((Nt)^{C_\ell}+C_\ell B\).  All these bounds include the requisite
finiteness assertions.  Set
\[
 D=\sum_{\ell\in\mathcal B_\chi}C_\ell,
 \qquad b=|\mathcal B_\chi|,
 \qquad C=C_0+D+b+bD.
\]
For \(x=Nt\ge2\), the total is at most
\[
 x^{C_0}+b x^D+bDB
 \le (b+1)x^{C_0+D}+CB
 \le x^{C_0+D+b}+CB
 \le x^C+CB,
\]
because \(b+1\le2^b\le x^b\).  This proves
\eqref{eq:host-point-bound}.  The source proves these steps in
\path{HeightAndVanishing/FinalLe.lean}, \texttt{sumIn};
\path{HeightAndVanishing/SumOut.lean}, \texttt{sumOut}; and
\path{HeightAndVanishing/PointArith.lean}, \texttt{point}.
The layer condition is \texttt{LayerLe} in
\path{HeightAndVanishing/LayerLe.lean};
\texttt{Lchi} in \path{PrimeSetBoundDataOfPrimitiveChar.lean}
is the full cyclotomic field of level \(\ell m\), and
\texttt{IsLayer} is in \path{StmtFields.lean}.

Every assembly declaration just cited, and the complete proofs of
\texttt{sumIn}, \texttt{sumOut}, \texttt{countIV}, and \texttt{point},
were read for this source reconstruction.  The imported arithmetic
and automorphic theorems are used under the same stipulated source
contract as \texttt{CorGK}; this is not a fresh build of the full
Lean project.  The explicit proof composition is recorded separately
as an uncompiled review fragment.
